\section{Model and Inference}
\label{sec:inference}

This section is standard machinery and carries no methodological claim.

\paragraph{Factors.} Each continuous factor uses one set of Fourier features built from the
routed mixture spectrum,
\begin{equation}
  s_{dr} = \sum_q \route_{drq} s_{dq},
  \qquad
  f_{dr}(x) = \phi_{s_{dr}}(x)^\top a_{dr},
  \qquad
  a_{dr} \sim \Norm(0, I),
  \label{eq:collapsed}
\end{equation}
so the variational coefficient budget is $2DR(1+2K)$, independent of the number of atoms in
the bank (\Cref{app:fairness}).

\paragraph{Variational family and objective.} Mean-field Gaussian
$q(a_{dr}) = \Norm(\mu_{dr}, \diag(\sigma_{dr}^2))$ over the Fourier coefficients; the CP core
$c_r$, routing logits $\alpha$ and observation noise $\sigma_y$ are point estimates.  Training
maximises
\begin{equation}
  \mathcal{F} = \Exp_q\Bigl[\textstyle\sum_{i\in\Obs}\log p(y_i \mid f, c, \sigma_y)\Bigr]
  - \sum_{d,r}\KL\bigl[q(a_{dr}) \,\|\, p(a_{dr})\bigr],
  \label{eq:elbo}
\end{equation}
with three reparameterised Monte-Carlo samples, Adam, and a fixed 400-step budget shared by
every method.  Numerical details are in \Cref{app:numerical}.

\paragraph{What the uncertainty is.} After training, $64$ posterior predictive samples at
$1024$ unobserved locations give $95\%$ interval coverage, Monte-Carlo predictive NLL, and
interval width, evaluated against independently re-noised held-out observations.  This is a
\emph{finite-Fourier, mean-field variational} posterior predictive, not an analytic exact GP
posterior, and $\alpha$ is an ELBO-optimised kernel hyperparameter rather than a Bayesian
model probability --- so the reported uncertainty excludes kernel-routing uncertainty.
